\documentclass[10pt,a4paper]{amsart}
\usepackage{amsmath,amssymb,amsthm,graphicx}
\usepackage{verbatim}
\newenvironment{reference}{\comment}{\endcomment}
%\newenvironment{reference}{}{}
\newenvironment{slogan}{\comment}{\endcomment}
\newenvironment{history}{\comment}{\endcomment}

% For commutative diagrams we use Xy-pic
\usepackage[all]{xy}

% We use 2cell for 2-commutative diagrams.
\xyoption{2cell}
% Theorem environments.
%
\theoremstyle{plain}
\newtheorem{theorem}[subsection]{Theorem}
\newtheorem{proposition}[subsection]{Proposition}
\newtheorem{lemma}[subsection]{Lemma}

\theoremstyle{definition}
\newtheorem{definition}[subsection]{Definition}
\newtheorem{example}[subsection]{Example}
\newtheorem{exercise}[subsection]{Exercise}
\newtheorem{situation}[subsection]{Situation}

\theoremstyle{remark}
\newtheorem{remark}[subsection]{Remark}
\newtheorem{remarks}[subsection]{Remarks}

\numberwithin{equation}{subsection}

% Macros
%
\def\lim{\mathop{\mathrm{lim}}\nolimits}
\def\colim{\mathop{\mathrm{colim}}\nolimits}
\def\Spec{\mathop{\mathrm{Spec}}}
\def\Hom{\mathop{\mathrm{Hom}}\nolimits}
\def\Ext{\mathop{\mathrm{Ext}}\nolimits}
\def\SheafHom{\mathop{\mathcal{H}\!\mathit{om}}\nolimits}
\def\SheafExt{\mathop{\mathcal{E}\!\mathit{xt}}\nolimits}
\def\Sch{\mathit{Sch}}
\def\Mor{\mathop{\mathrm{Mor}}\nolimits}
\def\Ob{\mathop{\mathrm{Ob}}\nolimits}
\def\Sh{\mathop{\mathit{Sh}}\nolimits}
\def\NL{\mathop{N\!L}\nolimits}
\def\CH{\mathop{\mathrm{CH}}\nolimits}
\def\proetale{{pro\text{-}\acute{e}tale}}
\def\etale{{\acute{e}tale}}
\def\QCoh{\mathit{QCoh}}
\def\Ker{\mathop{\mathrm{Ker}}}
\def\Im{\mathop{\mathrm{Im}}}
\def\Coker{\mathop{\mathrm{Coker}}}
\def\Coim{\mathop{\mathrm{Coim}}}

% Boxtimes
%
\DeclareMathSymbol{\boxtimes}{\mathbin}{AMSa}{"02}

%
% Macros for moduli stacks/spaces
%
\def\QCohstack{\mathcal{QC}\!\mathit{oh}}
\def\Cohstack{\mathcal{C}\!\mathit{oh}}
\def\Spacesstack{\mathcal{S}\!\mathit{paces}}
\def\Quotfunctor{\mathrm{Quot}}
\def\Hilbfunctor{\mathrm{Hilb}}
\def\Curvesstack{\mathcal{C}\!\mathit{urves}}
\def\Polarizedstack{\mathcal{P}\!\mathit{olarized}}
\def\Complexesstack{\mathcal{C}\!\mathit{omplexes}}
% \Pic is the operator that assigns to X its picard group, usage \Pic(X)
% \Picardstack_{X/B} denotes the Picard stack of X over B
% \Picardfunctor_{X/B} denotes the Picard functor of X over B
\def\Pic{\mathop{\mathrm{Pic}}\nolimits}
\def\Picardstack{\mathcal{P}\!\mathit{ic}}
\def\Picardfunctor{\mathrm{Pic}}
\def\Deformationcategory{\mathcal{D}\!\mathit{ef}}
\usepackage{cleveref}
\crefname{originallemma}{original lemma}{original lemmas}
\Crefname{originallemma}{Original lemma}{Original lemmas}
\makeatletter
\newlabel{algebra-lemma-support-point}{{00L5}{source}{}{sourcepointer-1}{}}
\newlabel{algebra-lemma-support-point@cref}{{[originallemma][0][]00L5}{[1][0][]source}{}{}{}}
\makeatother
\crefname{originallemma}{original lemma}{original lemmas}
\Crefname{originallemma}{Original lemma}{Original lemmas}
\makeatletter
\newlabel{coherent-lemma-coherent-Noetherian}{{01XZ}{source}{}{sourcepointer-2}{}}
\newlabel{coherent-lemma-coherent-Noetherian@cref}{{[originallemma][0][]01XZ}{[1][0][]source}{}{}{}}
\makeatother
\crefname{originallemma}{original lemma}{original lemmas}
\Crefname{originallemma}{Original lemma}{Original lemmas}
\makeatletter
\newlabel{decent-spaces-lemma-decent-no-specializations-map-to-same-point}{{03K5}{source}{}{sourcepointer-3}{}}
\newlabel{decent-spaces-lemma-decent-no-specializations-map-to-same-point@cref}{{[originallemma][0][]03K5}{[1][0][]source}{}{}{}}
\makeatother
\crefname{originallemma}{original lemma}{original lemmas}
\Crefname{originallemma}{Original lemma}{Original lemmas}
\makeatletter
\newlabel{decent-spaces-lemma-decent-specialization}{{03IL}{source}{}{sourcepointer-4}{}}
\newlabel{decent-spaces-lemma-decent-specialization@cref}{{[originallemma][0][]03IL}{[1][0][]source}{}{}{}}
\makeatother
\crefname{originallemma}{original lemma}{original lemmas}
\Crefname{originallemma}{Original lemma}{Original lemmas}
\makeatletter
\newlabel{decent-spaces-lemma-locally-Noetherian-decent-quasi-separated}{{0BB6}{source}{}{sourcepointer-5}{}}
\newlabel{decent-spaces-lemma-locally-Noetherian-decent-quasi-separated@cref}{{[originallemma][0][]0BB6}{[1][0][]source}{}{}{}}
\makeatother
\crefname{originalproposition}{original proposition}{original propositions}
\Crefname{originalproposition}{Original proposition}{Original propositions}
\makeatletter
\newlabel{decent-spaces-proposition-reasonable-sober}{{03K6}{source}{}{sourcepointer-6}{}}
\newlabel{decent-spaces-proposition-reasonable-sober@cref}{{[originalproposition][0][]03K6}{[1][0][]source}{}{}{}}
\makeatother
\crefname{originallemma}{original lemma}{original lemmas}
\Crefname{originallemma}{Original lemma}{Original lemmas}
\makeatletter
\newlabel{morphisms-lemma-support-finite-type}{{056J}{source}{}{sourcepointer-7}{}}
\newlabel{morphisms-lemma-support-finite-type@cref}{{[originallemma][0][]056J}{[1][0][]source}{}{}{}}
\makeatother
\crefname{originallemma}{original lemma}{original lemmas}
\Crefname{originallemma}{Original lemma}{Original lemmas}
\makeatletter
\newlabel{schemes-lemma-scheme-sober}{{01IS}{source}{}{sourcepointer-8}{}}
\newlabel{schemes-lemma-scheme-sober@cref}{{[originallemma][0][]01IS}{[1][0][]source}{}{}{}}
\makeatother
\crefname{originallemma}{original lemma}{original lemmas}
\Crefname{originallemma}{Original lemma}{Original lemmas}
\makeatletter
\newlabel{spaces-cohomology-lemma-coherent-Noetherian}{{07UB}{source}{}{sourcepointer-9}{}}
\newlabel{spaces-cohomology-lemma-coherent-Noetherian@cref}{{[originallemma][0][]07UB}{[1][0][]source}{}{}{}}
\makeatother
\crefname{originallemma}{original lemma}{original lemmas}
\Crefname{originallemma}{Original lemma}{Original lemmas}
\makeatletter
\newlabel{spaces-morphisms-lemma-support-finite-type}{{07TZ}{source}{}{sourcepointer-10}{}}
\newlabel{spaces-morphisms-lemma-support-finite-type@cref}{{[originallemma][0][]07TZ}{[1][0][]source}{}{}{}}
\makeatother
\title{Finite nonzero length at a point characterises the generic points of the support of a coherent sheaf}
\author{The Stacks Project Authors}
\date{stacks-project:spaces-chow.tex:local-snapshot-20261012; GNU FDL 1.2 or later}
\begin{document}
\maketitle
\noindent\textit{Editorial note.} All delivered mathematical text is the author's own, copied from the pinned Stacks Project chapter file spaces-chow.tex; no statement or proof step was rewritten, reordered or removed.\par\medskip
\noindent\textit{Editorial note.} The counted proof is the author's complete argument: an etale chart translation of support and length into a statement about the stalk, settled by an algebra criterion, followed by a specialisation and soberness argument for the decent case\par\medskip
\noindent\textit{Editorial note.} The label lemma-length-finite names the counted statement, and the proof environment immediately following it in the source is delivered verbatim as its proof.\par\medskip
\noindent\textit{Editorial note.} The counted statement is carried by the author's own lemma environment in the chapter source.\par\medskip
\noindent\textit{Allowed standard prerequisites (editorial).}\begin{itemize}
\item The length of a coherent module at a point of a locally Noetherian algebraic space and the support of a coherent module, together with the compatibility of both notions with etale base change to a scheme
\item The algebra criterion identifying the generic points of the support of a finite module over a Noetherian local ring with the points where the support is the maximal ideal, and finiteness of the module over such a ring
\item For the decent case: the lifting of specialisations along an etale morphism between spaces and the fact that the underlying spaces of schemes and of decent spaces are sober
\end{itemize}\medskip
\ifdefined\hypertarget\hypertarget{sourcepointer-1}{}\fi
\noindent\textit{Source reference.} Original lemma 00L5: sources/stacks-chow/tags.txt, source lines 775--775. Prerequisite pointer invoked at proof line(s) 599; it names a standard fact of a sibling chapter that the counted argument uses, and every step of the counted proof built on it is retained verbatim in the delivered body.\par\medskip
\ifdefined\hypertarget\hypertarget{sourcepointer-2}{}\fi
\noindent\textit{Source reference.} Original lemma 01XZ: sources/stacks-chow/tags.txt, source lines 2538--2538. Prerequisite pointer invoked at proof line(s) 589; it names a standard fact of a sibling chapter that the counted argument uses, and every step of the counted proof built on it is retained verbatim in the delivered body.\par\medskip
\ifdefined\hypertarget\hypertarget{sourcepointer-3}{}\fi
\noindent\textit{Source reference.} Original lemma 03K5: sources/stacks-chow/tags.txt, source lines 4665--4665. Prerequisite pointer invoked at proof line(s) 611; it names a standard fact of a sibling chapter that the counted argument uses, and every step of the counted proof built on it is retained verbatim in the delivered body.\par\medskip
\ifdefined\hypertarget\hypertarget{sourcepointer-4}{}\fi
\noindent\textit{Source reference.} Original lemma 03IL: sources/stacks-chow/tags.txt, source lines 4607--4607. Prerequisite pointer invoked at proof line(s) 610; it names a standard fact of a sibling chapter that the counted argument uses, and every step of the counted proof built on it is retained verbatim in the delivered body.\par\medskip
\ifdefined\hypertarget\hypertarget{sourcepointer-5}{}\fi
\noindent\textit{Source reference.} Original lemma 0BB6: sources/stacks-chow/tags.txt, source lines 14234--14234. Prerequisite pointer invoked at proof line(s) 622; it names a standard fact of a sibling chapter that the counted argument uses, and every step of the counted proof built on it is retained verbatim in the delivered body.\par\medskip
\ifdefined\hypertarget\hypertarget{sourcepointer-6}{}\fi
\noindent\textit{Source reference.} Original proposition 03K6: sources/stacks-chow/tags.txt, source lines 4666--4666. Prerequisite pointer invoked at proof line(s) 613; it names a standard fact of a sibling chapter that the counted argument uses, and every step of the counted proof built on it is retained verbatim in the delivered body.\par\medskip
\ifdefined\hypertarget\hypertarget{sourcepointer-7}{}\fi
\noindent\textit{Source reference.} Original lemma 056J: sources/stacks-chow/tags.txt, source lines 6712--6712. Prerequisite pointer invoked at proof line(s) 591; it names a standard fact of a sibling chapter that the counted argument uses, and every step of the counted proof built on it is retained verbatim in the delivered body.\par\medskip
\ifdefined\hypertarget\hypertarget{sourcepointer-8}{}\fi
\noindent\textit{Source reference.} Original lemma 01IS: sources/stacks-chow/tags.txt, source lines 1989--1989. Prerequisite pointer invoked at proof line(s) 615; it names a standard fact of a sibling chapter that the counted argument uses, and every step of the counted proof built on it is retained verbatim in the delivered body.\par\medskip
\ifdefined\hypertarget\hypertarget{sourcepointer-9}{}\fi
\noindent\textit{Source reference.} Original lemma 07UB: sources/stacks-chow/tags.txt, source lines 9966--9966. Prerequisite pointer invoked at proof line(s) 587; it names a standard fact of a sibling chapter that the counted argument uses, and every step of the counted proof built on it is retained verbatim in the delivered body.\par\medskip
\ifdefined\hypertarget\hypertarget{sourcepointer-10}{}\fi
\noindent\textit{Source reference.} Original lemma 07TZ: sources/stacks-chow/tags.txt, source lines 9954--9954. Prerequisite pointer invoked at proof line(s) 605; it names a standard fact of a sibling chapter that the counted argument uses, and every step of the counted proof built on it is retained verbatim in the delivered body.\par\medskip
\begin{lemma}
\label{lemma-length-finite}
Let $S$ be a scheme and let $X$ be a
locally Noetherian algebraic space over $S$.
Let $\mathcal{F}$ be a coherent $\mathcal{O}_X$-module.
Let $x \in |X|$. The following are equivalent
\begin{enumerate}
\item for some \'etale morphism $U \to X$ with $U$ a scheme
and $u \in U$ mapping to $x$ we have $u$ is a generic point
of an irreducible component of $\text{Supp}(\mathcal{F}|_U)$,
\item for any \'etale morphism $U \to X$ with $U$ a scheme
and $u \in U$ mapping to $x$ we have $u$ is a generic point
of an irreducible component of $\text{Supp}(\mathcal{F}|_U)$,
\item the length of $\mathcal{F}$ at $x$ is finite and nonzero.
\end{enumerate}
If $X$ is decent (equivalently quasi-separated) then these are
also equivalent to
\begin{enumerate}
\item[(4)] $x$ is a generic point of an irreducible component of
$\text{Supp}(\mathcal{F})$.
\end{enumerate}
\end{lemma}

\begin{proof}
Assume $f : U \to X$ is an \'etale morphism with $U$ a scheme
and $u \in U$ maps to $x$. Then $\mathcal{F}|_U = f^*\mathcal{F}$
is a coherent $\mathcal{O}_U$-module on the locally Noetherian
scheme $U$ and in particular $(\mathcal{F}|_U)_u$ is a finite
$\mathcal{O}_{U, u}$-module, see Cohomology of Spaces, Lemma
\ref{spaces-cohomology-lemma-coherent-Noetherian}
and Cohomology of Schemes, Lemma
\ref{coherent-lemma-coherent-Noetherian}.
Recall that the support of $\mathcal{F}|_U$ is a closed subset of $U$
(Morphisms, Lemma \ref{morphisms-lemma-support-finite-type}) and
that the support of $(\mathcal{F}|_U)_u$ is the pullback
of the support of $\mathcal{F}|_U$ by the morphism
$\Spec(\mathcal{O}_{U, u}) \to U$. Thus $u$ is a generic point
of an irreducible component of $\text{Supp}(\mathcal{F}|_U)$
if and only if the support of $(\mathcal{F}|_U)_u$ is equal
to the maximal ideal of $\mathcal{O}_{U, u}$.
Now the equivalence of (1), (2), (3) follows from
by Algebra, Lemma \ref{algebra-lemma-support-point}.

\medskip\noindent
If $X$ is decent we choose an \'etale morphism $f : U \to X$ and a point
$u \in U$ mapping to $x$. The support of $\mathcal{F}$ pulls back to
the support of $\mathcal{F}|_U$, see
Morphisms of Spaces, Lemma \ref{spaces-morphisms-lemma-support-finite-type}.
Also, specializations $x' \leadsto x$ in $|X|$ lift to specializations
$u' \leadsto u$ in $U$ and any nontrivial specialization $u' \leadsto u$ in $U$
maps to a nontrivial specialization $f(u') \leadsto f(u)$ in $|X|$, see
Decent Spaces, Lemmas
\ref{decent-spaces-lemma-decent-specialization} and
\ref{decent-spaces-lemma-decent-no-specializations-map-to-same-point}.
Using that $|X|$ and $U$ are sober topological spaces
(Decent Spaces, Proposition \ref{decent-spaces-proposition-reasonable-sober}
and
Schemes, Lemma \ref{schemes-lemma-scheme-sober}) we conclude
$x$ is a generic point of the support of $\mathcal{F}$
if and only if $u$ is a generic point of the support of
$\mathcal{F}|_U$. We conclude (4) is equivalent to (1).

\medskip\noindent
The parenthetical statement follows from Decent Spaces, Lemma
\ref{decent-spaces-lemma-locally-Noetherian-decent-quasi-separated}.
\end{proof}

\end{document}
